\chapter{Remaining Work}
\label{chap:future-work}

This chapter states plainly what has not been built yet, so that nothing in
the preceding chapters is mistaken for a finished system. Everything listed
here is planned, in the project's own staged schedule, for implementation
steps I6--I7 and evaluation steps e0--e7; none of it exists in the
repository at the time of writing, confirmed by direct inspection rather
than assumed.

\section{Remaining Implementation}

\paragraph{Full pipeline orchestration (I7).} \Cref{chap:architecture}
describes an orchestrator that loops over every usable row, running
classification, explanation, repair, and fix application in sequence and
logging every attempt, including ones where a component itself fails. No
such loop exists yet: each component so far has been run and validated
independently, as described in \Cref{chap:methodology,chap:validation}.
Building it is mostly integration work, since every component it would call
already has a stable, tested interface.

\paragraph{Bounded two-round repair (ablation e3).} When FixApplicator
reports \texttt{still\_failing} with a genuinely new error
(\Cref{sec:v-key-observation}), the design allows feeding that new error
back into RAGRepairAgent for one further attempt. The pilot in
\Cref{chap:validation} shows this input already exists and is already
correctly distinguished from "the fix did not work"; the feedback loop
itself, and the planned one-round-versus-two-round comparison, have not
been built or run.

\paragraph{ResultLogger and \texttt{repair\_attempts} persistence (O4).} No
component yet writes to the \texttt{repair\_attempts} table designed in
\Cref{subsec:repair-table}. Every result produced so far is written to a
plain JSONL file by the component that produced it. Building ResultLogger
means creating the SQLite table and writing one row per attempt from the
classification, explanation, fix, and outcome already produced by the
existing components.

\paragraph{Knowledge-graph enrichment (I6, O6).} The RML mapping described
in \Cref{subsec:kg} that would turn a \texttt{repair\_attempts} row into RDF
triples inside the FAIR Jupyter knowledge graph has not been written. This
step depends on ResultLogger existing first, since it maps from that
table.

\section{Remaining Evaluation}

\paragraph{Full-dataset evaluation run (e1--e2).} The pilots in
\Cref{chap:validation} cover fewer than ten records in total. The 187-row
evaluation split reserved in \Cref{chap:problem-dataset} has not been run.
Doing so requires the orchestrator above, since running 187 rows one at a
time by hand is not practical.

\paragraph{One-round versus two-round ablation (e3).} Comparing repair
success with and without the second repair round requires both the
orchestrator and the two-round mode described above; neither exists yet.

\paragraph{User study and inter-rater agreement (e4--e5).} The project plan
includes a user study with Likert-scale scoring of explanation quality and
a Cohen's kappa analysis of agreement between raters. This has not started;
it depends on having a stable set of generated explanations to show
participants, which \Cref{chap:validation} shows exists only for a handful
of pilot rows so far.

\paragraph{Baseline comparison (e6).} Comparing this system's repair success
rate against a baseline (for example, a version of the repair agent without
PyPI grounding) requires the full evaluation run above to exist first, so
the two can be compared on the same rows.

\paragraph{Benchmark dataset finalisation (e7).} Objective O4 calls for
publishing the completed \texttt{repair\_attempts} table as a reusable
benchmark. This is the last step in the chain above, since it depends on
every step before it.

\section{Summary}

\Cref{tab:future-summary} separates what this thesis reports as done from
what remains, to avoid any ambiguity.

\begin{table}[h]
  \centering
  \begin{tabular}{|p{7.2cm}|l|}
    \hline
    \textbf{Item} & \textbf{Status} \\
    \hline
    Dataset preparation, classification, explanation, PyPI-grounded repair
      proposal generation, fix application and re-execution (I1--I5) &
      Implemented, tested, pilot-validated \\
    \hline
    Full pipeline orchestrator (I7) & Not started \\
    \hline
    Two-round repair loop (e3) & Not started \\
    \hline
    ResultLogger / \texttt{repair\_attempts} persistence (O4) & Not started \\
    \hline
    Knowledge-graph enrichment (I6, O6) & Not started \\
    \hline
    Full 187-row evaluation run (e1--e2) & Not started \\
    \hline
    One-round vs.\ two-round ablation (e3) & Not started \\
    \hline
    User study and Cohen's kappa analysis (e4--e5) & Not started \\
    \hline
    Baseline comparison (e6) & Not started \\
    \hline
    Benchmark dataset publication (e7) & Not started \\
    \hline
  \end{tabular}
  \caption{Completed work versus remaining work, at the time of writing.}
  \label{tab:future-summary}
\end{table}
